\subsection{NN-G030: associativity forces modulo-four constraints beyond trace parity}
\paragraph{Statement.}
For every integral associative closed rank-three table with
the specified trace parities, put $\sigma=D(D-1)$ and
$A_2=\delta_2\Lambda_2$, the actual second cell-product
coefficient. No constant-ratio, degree, or auxiliary-integrality
hypothesis is needed. Then the following are polynomial
congruences in $\mathbb Z[D]$:
\[
 S\equiv\sigma G_1,\qquad
 G_2-D\delta_2\equiv G_1,\qquad
 DG_2-A_2\equiv G_1\pmod4.                           \tag{G030.1}
\]
In fact
\[
 \begin{aligned}
 G_1&\in(4,2D),\\
 G_2,A_2&\in(4,2D,D^2),\\
 S&\in(4,2D^2(D-1)).
 \end{aligned}                                                       \tag{G030.2}
\]
Thus $G_1,G_2,S,A_2$ take values divisible by four
at every even integer $D$, and $S$ does so at every
integer $D$. In particular their ordinary dimension
masses at $D=4$ are multiples of four.
These necessary conditions use associativity and
integral trace parity alone, before effectivity or
domination is tested.

\paragraph{Proof.}
The supplied trace table gives
\[
 \begin{gathered}
 W\equiv D,\quad\delta_2\equiv\sigma,\quad
 G_1\equiv S\equiv0,\quad G_2\equiv D\sigma,\\
 U_2\equiv V_2\equiv D,\qquad A_2\equiv D^2\sigma
       \pmod2.
 \end{gathered}
\]
Write $G_1=2H$ and $S=2K$, and define the integral
polynomials
\[
 P=(G_2-D\delta_2)/2,\qquad
 Q=(DG_2-A_2)/2.
\]
The three necessary associativity identities reduce to
\[
 WK=HG_2,\qquad U_2K=G_2P,\qquad V_2K=G_2Q.
\]
Reducing the first modulo two and cancelling the
nonzero polynomial $D$ in $\mathbb F_2[D]$ gives
$\overline K=\overline H\sigma$.
The second and third then give
\[
 D\overline H\sigma=D\sigma\overline P
                       =D\sigma\overline Q.
\]
Cancelling $D\sigma\ne0$ gives
$\overline P=\overline Q=\overline H$.
Multiplying these statements by two proves (G030.1).
This does not require $\overline H$ itself to be nonzero.

At $D=0$, (G030.1) says
\[
 S(0)\equiv0,\quad G_2(0)\equiv G_1(0),\quad
 A_2(0)\equiv-G_1(0)\pmod4.
\]
The integers $G_1(0),G_2(0)$ are even, and $W(0)$
is even. Since $WS=G_1G_2$, their product at zero
is divisible by eight. If their common residue
modulo four were two, their product would instead
be four modulo eight. Hence $G_1(0),G_2(0),A_2(0)$
are all divisible by four.

Every coefficient of $G_1$ is even, so its divisible-by-four
constant term gives $G_1\in(4,2D)$. The parities
$G_2\equiv D^2(D-1)$ and
$A_2\equiv D^3(D-1)$ show that their linear
coefficients are even; their constant terms are
multiples of four. This gives their ideal membership
in (G030.2). Finally substitute $G_1=4a+2DB$
in the first congruence of (G030.1):
\[
 S\equiv2D^2(D-1)B\pmod4,
\]
which proves the last ideal membership. Evaluation
at even integers gives the claimed divisibilities.
For arbitrary integer $D$, the product $D^2(D-1)$
is even, proving the additional assertion for $S$.
\hfill$\square$

\paragraph{A low-mass scalar candidate that these constraints exclude.}
Ordinary mass parity and the scalar associativity
identities alone admit the following hypothetical data:
\[
 \begin{array}{c|rrrrrrrrrrrr}
 &w&g_1&g_2&s&d_1&d_2&u_1&u_2&v_1&v_2&a_1&a_2\\\hline
 \text{values}&10&70&90&630&4&12&6&6&30&48&10&24
 \end{array}
\]
Here lower-case symbols denote evaluations at $D=4$,
with $a_i=A_i(4)$. They obey all target mass lower
bounds and are even. Directly,
\[
 ws=g_1g_2,
 \qquad u_i s=g_i(g_i-4d_i),\qquad
 v_i s=g_i(4g_i-a_i),
\]
For example both $u_i s$ equal $3780$;
$v_1s=18900$ and $v_2s=30240$.
Their relaxed total mass would be $8522$, or
$(8522-4424)/2=2049$ pairs. But $g_1=70$, $g_2=90$
and $s=630$ are not multiples of four. Thus no integral
polynomial table with the required trace parity can
realize these scalar data. This is a strict exclusion
from the scalar relaxation, not an assertion that the
hypothetical data were an actual lower-cost lift.

The constraints do not say that every coefficient
of $G_1$ or $S$ is divisible by four. In the actual
$3352$ attainer the coefficient of $D$ in $G_1$
is six and the coefficient of $D^2$ in $S$ is
$-94$. Both are two modulo four. The polynomial
ideals in (G030.2), rather than coefficientwise
divisibility by four, are the proved statements.
The small exact checker verifies the displayed
congruences on that attainer and verifies all scalar
identities in the excluded mass example.

\paragraph{Geometric translation and algorithmic scope.}
At the level of integer congruences, individual effective
graded sign representations require dimension and signed
trace to agree modulo two.
Associativity of their finite-kernel composition imposes
the additional congruences (G030.1). In the symmetric
character ring $\mathbb Z[D]$, they say that the listed
coefficient characters vanish in the explicit quotient
rings $\mathbb Z[D]/(4,2D)$,
$\mathbb Z[D]/(4,2D,D^2)$, and
$\mathbb Z[D]/(4,2D^2(D-1))$, as appropriate.
These are exact scheme-theoretic congruence statements
over the nonreduced characteristic-four base, not
purity or point-counting assertions.

The character interpretation is SGA 3, I, 4.7.2--4.7.3;
$D=4$ is evaluation at the identity grading character,
while $D=0$ is evaluation at $q=-1$.
The composing objects remain the NN-G010 kernel
realization of archive cote 2, scan 17. These conditions
provide inexpensive necessary tests before a bounded
search or NP verification, but are not sufficient for
an effective lift and do not determine its optimum,
canonicality, rank-four positivity, or PH0/PH1.
